% GENERATED FILE. Edit canonical lessons, then run npm run generate:books.
% canonical-source-hash: fnv1a64:b3160f4ad56aa62d
% canonical-lessons: IT-C29-si, IT-C29-no, IT-C29-si-no-grazie, IT-C29-daccordo, IT-C29-vero

\chapter{Yes and No}
\label{ch:it-si-e-no}
\hlchaptermodality{29}

\begin{chapteropening}
\textbf{By the end of this chapter:} \emph{I can say yes and no, take or turn down what is offered, agree, and ask somebody to confirm what I have just said.}

The chapter's last lesson gives vero its second job — the tag hung on the end of a statement to ask the other person to confirm it — and gathers the four answering moves the chapter opened: sì, no, sì grazie / no grazie, and d'accordo.
\end{chapteropening}

\section[sì]{sì — yes, and where a language gets one}

\label{lesson:IT-C29-si}



\begin{quote}
A new chapter, so the recalls come first, at the four growing distances this book measures.

\begin{itemize}
\raggedright
  \item \emph{Recall:} say \emph{prima il primo, dopo il secondo} — \textbf{R1}, one lesson back, before it decays
  \item \emph{Recall:} say \emph{decimo / decima} — \textbf{R2}, the first real retrieval, five to fifteen lessons back
  \item \emph{Recall:} say \emph{scrivere} — \textbf{R3}, durable, twenty to sixty lessons back
  \item \emph{Recall:} say \emph{prendere} — \textbf{R3}, durable again, a second word from that distance
  \item \emph{Recall:} say \emph{arrivederci} — \textbf{R4}, recognition at distance, eighty lessons back or more
\end{itemize}
\end{quote}

\subsection*{You'll want to know: The word}
\begin{quote}
\textbf{sì} — \textbf{yes}.
\end{quote}

It has been in front of you since the first \emph{-are} verbs. That chapter's dialogue ran \emph{— Parli italiano? — \textbf{Sì}, parlo un po'.} The word was said, and it was never given a line of its own, because the lesson was about the \emph{-are} endings. Here it is on its own.

The accent is not decoration. It is written \textbf{sì}, with a grave accent, to keep it apart from \textbf{si}, an unstressed little word with several other jobs — one of which you have already met inside \emph{così così}.

\begin{cousinweb}
Classical Latin had \textbf{no word for yes}. Asked \emph{Venīsne?} — \textquotedblleft{}Are you coming?\textquotedblright{} — a Roman answered by repeating the verb: \emph{veniō}, \textquotedblleft{}I am coming.\textquotedblright{} Every Romance language had to build a \emph{yes} out of something else, and each one chose differently. French took \emph{hoc ille} $\to$ \emph{oui}. Italian, Spanish and Catalan took \textbf{sīc}, \textquotedblleft{}so, thus\textquotedblright{} — the answer meaning \emph{that is how it is}.

You already own that \emph{sīc}. The wellbeing chapter gave you \textbf{così così}, \textquotedblleft{}so-so\textquotedblright{}, and \emph{così} is \emph{eccu(m) sīc}, \textquotedblleft{}behold, thus\textquotedblright{}. \textbf{Sì} and \emph{così} are the same Latin word, one of them worn down to a single syllable.
\end{cousinweb}

\subsection*{Your turn}
Say these aloud:

\begin{itemize}
\raggedright
  \item \emph{sì}
  \item \emph{Parli italiano? — Sì.}
  \item \emph{sì} and \emph{così} — one Latin word, two Italian ones
\end{itemize}

\begin{tcolorbox}[breakable,colback=teal!4,colframe=teal!35!black,title={Before you move on}]
How do you say yes? (\textbf{Sì}.) Which Latin word is it? (\emph{\textbf{Sīc}}, \textquotedblleft{}so, thus\textquotedblright{}.) Which word already in this book is built on the same \emph{sīc}? (\emph{\textbf{Così}}, in \emph{così così}.)
\end{tcolorbox}
\section[no]{no — the other half of the answer}

\label{lesson:IT-C29-no}



\begin{quote}
Two recalls from this chapter and this book's four measured distances.

\begin{itemize}
\raggedright
  \item \emph{Recall:} say \emph{sì} — \textbf{R1}, one lesson back, before it decays
  \item \emph{Recall:} say \emph{undicesimo} — \textbf{R2}, the first real retrieval, five to fifteen lessons back
  \item \emph{Recall:} say \emph{chiedere} — \textbf{R3}, durable, twenty to sixty lessons back
  \item \emph{Recall:} say \emph{aiutare} — \textbf{R3}, durable again, a second word from that distance
  \item \emph{Recall:} say \emph{a domani} — \textbf{R4}, recognition at distance, eighty lessons back or more
\end{itemize}
\end{quote}

\subsection*{You'll want to know: The word}
\begin{quote}
\textbf{no} — \textbf{no}.
\end{quote}

\emph{Sì} and \textbf{no} are the pair. One of them answers a question by agreeing with it; the other answers by rejecting it, and it can stand entirely alone:

\begin{quote}
— \textbf{Parli italiano?} — \textbf{No.}
\end{quote}

That is a complete Italian turn. Nothing else is required of you.

\begin{cousinweb}
\textbf{No} is Latin \textbf{nōn}. So is the word in the next lesson. Italian took a single Latin negator and split it by stress into two words that no longer look interchangeable:

\noindent
\begin{tabularx}{\linewidth}{@{}>{\raggedright\arraybackslash}X>{\raggedright\arraybackslash}X>{\raggedright\arraybackslash}X@{}}
\toprule
\textbf{} & \textbf{} & \textbf{} \\
\midrule
\textbf{no} & stressed, stands alone & \emph{— Parli italiano? — \textbf{No}.} \\
\textbf{non} & unstressed, leans on a verb & \emph{(next lesson)} \\
\bottomrule
\end{tabularx}

English does the same thing without noticing: \emph{no} and \emph{not} are also one old word pulled apart. Say \textbf{no} on its own, and never in front of a verb.
\end{cousinweb}

\subsection*{Your turn}
Say these aloud:

\begin{itemize}
\raggedright
  \item \emph{no}
  \item \emph{— Parli italiano? — No.}
  \item \emph{sì} … \emph{no} — the pair
\end{itemize}

\begin{tcolorbox}[breakable,colback=teal!4,colframe=teal!35!black,title={Before you move on}]
How do you say no? (\textbf{No}.) Which Latin word is behind it? (\emph{\textbf{Nōn}}.) And where may \emph{no} not stand? (\textbf{In front of a verb} — it answers on its own.)
\end{tcolorbox}
\section[sì, grazie / no, grazie]{sì, grazie — no, grazie}

\label{lesson:IT-C29-si-no-grazie}



\begin{quote}
Two words you were given days apart now do one job together.

\begin{itemize}
\raggedright
  \item \emph{Recall:} say \emph{no} — \textbf{R1}, one lesson back, before it decays
  \item \emph{Recall:} say \emph{prima} — \textbf{R2}, the first real retrieval, five to fifteen lessons back
  \item \emph{Recall:} say \emph{mi piace} — \textbf{R3}, durable, twenty to sixty lessons back
  \item \emph{Recall:} say \emph{portare} — \textbf{R3}, durable again, a second word from that distance
  \item \emph{Recall:} say \emph{a presto} — \textbf{R4}, recognition at distance, eighty lessons back or more
\end{itemize}
\end{quote}

\subsection*{You'll want to know: accepting and declining}
Somebody holds out a coffee. In Italian you need two words, and you have had one of them since the very first chapter.

\begin{quote}
\textbf{Sì, grazie.} — Yes, please. \emph{(taking it)} \textbf{No, grazie.} — No, thank you. \emph{(turning it down)}
\end{quote}

Note what Italian does \textbf{not} say. English says \textquotedblleft{}yes, \emph{please}\textquotedblright{} when it accepts and \textquotedblleft{}no, \emph{thank you}\textquotedblright{} when it declines: two different words. Italian uses \textbf{grazie} for both. The \emph{sì} or the \emph{no} carries the whole decision, and \emph{grazie} only softens it.

This is why \emph{grazie} alone is risky at an Italian table. Offered a second helping, a bare \textbf{grazie} with a small backward tilt of the head is heard as a refusal. If you want it, say \textbf{sì, grazie}.

\subsection*{Your turn}
Say these aloud:

\begin{itemize}
\raggedright
  \item \emph{Sì, grazie.}
  \item \emph{No, grazie.}
  \item \emph{— Un caffè? — Sì, grazie.}
\end{itemize}

\begin{tcolorbox}[breakable,colback=teal!4,colframe=teal!35!black,title={Before you move on}]
Somebody offers you a coffee and you want it — what do you say? (\textbf{Sì, grazie}.) And if you do not? (\textbf{No, grazie}.) Which word is the same in both? (\emph{\textbf{Grazie}} — Italian does not swap it for a \emph{please}.)
\end{tcolorbox}
\section[d'accordo]{d'accordo — agreed}

\label{lesson:IT-C29-daccordo}



\begin{quote}
Two chapters' worth of distance before a word about agreeing.

\begin{itemize}
\raggedright
  \item \emph{Recall:} say \emph{sì, grazie / no, grazie} — \textbf{R1}, one lesson back, before it decays
  \item \emph{Recall:} say \emph{dopo} — \textbf{R2}, the first real retrieval, five to fifteen lessons back
  \item \emph{Recall:} say \emph{comprare} — \textbf{R3}, durable, twenty to sixty lessons back
  \item \emph{Recall:} say \emph{aspettare} — \textbf{R3}, durable again, a second word from that distance
  \item \emph{Recall:} say \emph{a più tardi} — \textbf{R4}, recognition at distance, eighty lessons back or more
\end{itemize}
\end{quote}

\subsection*{You'll want to know: The phrase}
\begin{quote}
\textbf{D'accordo.} — \textbf{Agreed. All right. Okay.}
\end{quote}

It is the whole answer. Somebody proposes something and you accept the proposal rather than the object:

\begin{quote}
— \textbf{A domani?} — \textbf{D'accordo.}
\end{quote}

You can also put a person in front of it with \emph{essere}, which the \emph{essere} lesson gave you:

\begin{quote}
\textbf{Sono d'accordo.} — I agree.
\end{quote}

The apostrophe is the \emph{di} you already meet inside \emph{piacere \textbf{di} conoscerti}, with its vowel knocked out before another vowel — the same elision as \emph{l'acqua} and \emph{l'amico}.

\begin{cousinweb}
\textbf{Accordo} is Latin \textbf{ad} + \textbf{cor, cordis} — \textquotedblleft{}\textbf{heart}.\textquotedblright{} To reach an \emph{accordo} is literally to bring two hearts to the same place.

You were taught \emph{cordis} in the body-words chapter without being told this. \textbf{Il cuore}, the heart, is that same \emph{cor}. So is English \textbf{accord}, \textbf{record} (to bring back to the heart, which is where the Romans thought memory lived), \textbf{concord}, \textbf{discord} and \textbf{courage}.

And an \emph{accordo} in music is a chord: notes that agree.
\end{cousinweb}

\subsection*{Your turn}
Say these aloud:

\begin{itemize}
\raggedright
  \item \emph{d'accordo}
  \item \emph{Sono d'accordo.}
  \item \emph{— A domani? — D'accordo.}
\end{itemize}

\begin{tcolorbox}[breakable,colback=teal!4,colframe=teal!35!black,title={Before you move on}]
How do you say \textquotedblleft{}agreed\textquotedblright{}? (\textbf{D'accordo}.) Which body word from the body-words chapter is hiding inside it? (\emph{\textbf{Cuore}} — Latin \emph{cor, cordis}.) And what is the apostrophe standing in for? (\emph{\textbf{Di}}, with its vowel dropped before a vowel.)
\end{tcolorbox}
\section[vero?]{vero — true, and the question you hang on the end}

\label{lesson:IT-C29-vero}



\begin{quote}
The last word of the chapter, and five retrievals before it.

\begin{itemize}
\raggedright
  \item \emph{Recall:} say \emph{d'accordo} — \textbf{R1}, one lesson back, before it decays
  \item \emph{Recall:} say \emph{prima il primo, dopo il secondo} — \textbf{R2}, the first real retrieval, five to fifteen lessons back
  \item \emph{Recall:} say \emph{incontrare} — \textbf{R3}, durable, twenty to sixty lessons back
  \item \emph{Recall:} say \emph{giocare, suonare} — \textbf{R3}, durable again, a second word from that distance
  \item \emph{Recall:} say \emph{(practice)} — \textbf{R4}, recognition at distance, eighty lessons back or more
\end{itemize}
\end{quote}

\subsection*{You'll want to know: The word}
\begin{quote}
\textbf{vero} (masculine) · \textbf{vera} (feminine) — \textbf{true}.
\end{quote}

It behaves as any describing word does and changes its ending with what it describes, exactly as \emph{primo / prima} did.

\begin{grammarlens}[title={Grammar lens: the turn it does at the end of a sentence}]
Italian has no separate machinery for the little question English tacks on the end — \emph{isn't it? doesn't he? aren't you?} Italian says one word for all of them:

\begin{quote}
\textbf{Parli italiano, vero?} — You speak Italian, don't you? \textbf{Marco abita a Roma, vero?} — Marco lives in Rome, doesn't he?
\end{quote}

It asks the other person to confirm what you have just said, and the answer it expects is \textbf{sì}. Its opposite number is \textbf{no?}, which is blunter and expects the same \emph{sì}: \emph{Parli italiano, no?}

English speakers find this restful. \emph{Vero?} keeps one shape for every sentence.
\end{grammarlens}

\subsection*{Your turn}
Say these aloud:

\begin{itemize}
\raggedright
  \item \emph{vero, vera}
  \item \emph{Parli italiano, vero?}
  \item \emph{Marco abita a Roma, vero?}
\end{itemize}

\begin{tcolorbox}[breakable,colback=teal!4,colframe=teal!35!black,title={Before you move on}]
What does \emph{vero} mean on its own? (\textbf{True}.) What does it do at the end of a sentence? (\textbf{It asks the other person to confirm} — English \emph{isn't it, doesn't he, aren't you}.) Which answer does it expect? (\emph{\textbf{Sì}}.)
\end{tcolorbox}
